\begin{oxlemma}{Progress}{lemma:progress-app}
  If
  \oxtypjudge{\oxglobalctx}{\oxemptyctx}{\oxkontctx}{\oxvarctx}{\oxexpr}
  {\oxsitype}{\oxvarctx^\prime} and
  \oxstorevalidity{\oxglobalctx}{\oxvarctx}{\oxstore}, then either
  $\oxexpr$ is a value, $\oxexpr$ is an \oxabort{\,\dots\,}, or $\oxexists
  \oxstore^\prime, \oxexpr^\prime. \;
  \oxreduce{\oxglobalctx}{\oxstore}{\oxexpr}{\oxstore^\prime}{\oxexpr^\prime}$.
\end{oxlemma}

\paragraph{Proof} We proceed by induction on the derivation
\oxtypjudge{\oxglobalctx}{\oxemptyctx}{\oxkontctx}{\oxvarctx}{\oxexpr}
{\oxtype}{\oxvarctx^\prime}. \\

\oxprogresscase{T-Move}{\TMove}{\EMove}{
  Applying \Lemma{lemma:place-exprs-reduce} to
  $\oxmusafety{\oxtvarctx}{\oxvarctx}{\oxmut}{\oxplace}{\oxset{
      \oxloanpkg{\oxmut}{\oxplace}
  }}$, $\oxcomputetynoprov{\oxtvarctx}{\oxvarctx}{\oxmut}{\oxplace}{\oxsitype}$
  (from $\oxlookup{\oxvarctx}{\oxplace}{\oxsitype}$ by \oxname{TC-Place}), and
  $\oxstorevalidity{\oxglobalctx}{\oxvarctx}{\oxstore}$ to conclude that
  $\oxnorm{\oxstore}{\oxplace}{\_}{\_}{\oxvalue}$. Thus, we can step
  with \oxname{E-Move}.
}

\oxprogresscase{T-Copy}{\TCopy}{\ECopy}{
  Applying \Lemma{lemma:place-exprs-reduce} to
  $\oxmusafety{\oxtvarctx}{\oxvarctx}{\oximm}{\oxplaceexpr}{\oxset{
      \overline{\oxloan}
  }}$, $\oxcomputetynoprov{\oxtvarctx}{\oxvarctx}{\oximm}{\oxplaceexpr}
  {\oxsitype}$, and $\oxstorevalidity{\oxglobalctx}{\oxvarctx}{\oxstore}$ to
  conclude that $\oxnorm{\oxstore}{\oxplaceexpr}{\_}{\_}{\oxvalue}$. Thus, we can
  step with \oxname{E-Copy}.
}

\oxprogresscasewide{T-Borrow}{\TBorrow}{\EBorrow}{
  Applying \Lemma{lemma:place-exprs-reduce} to
  $\oxmusafety{\oxtvarctx}{\oxvarctx}{\oxmuta}{\oxplaceexpr}{\oxset{
      \overline{\oxloan}
    }}$, $\oxcomputetynoprov{\oxtvarctx}{\oxvarctx}{\oxmuta}{\oxplaceexpr}
  {\oxxitype}$, and $\oxstorevalidity{\oxglobalctx}{\oxvarctx}{\oxstore}$ to
  conclude that $\oxnorm{\oxstore}{\oxplaceexpr}{\oxreferent}{\_}{\_}$. Thus, we can
  step with \oxname{E-Borrow}.
}

\oxprogresscaseinductive{T-BorrowIndex}{\TBorrowIndex}{
  We proceed based on  whether or not $\oxexpr$ is a value. If it is not, we can
  decompose our expression into the evaluation context
  $\oxref{\oxprov}{\oxmuta}{\oxindex{\oxplaceexpr}{\oxhole}}$ and redex
  $\oxexpr$. Then, by applying our induction hypothesis to the typing derivation
  for $\oxexpr$, we know either that $\oxexpr$ is an $\oxkey{abort!}$ expression
  or it $\oxexpr$ steps to some $\oxexpr^\prime$.  In the former case, we can
  step with \oxname{E-EvalCtxAbort}.  In the latter case, we can plug
  $\oxexpr^\prime$ back into our evaluation context and step with
  \oxname{E-EvalCtx}.
}{
  If $\oxexpr$ is a value, we would like to step with one of:
}{\EBorrowIndex \and \EBorrowIndexOOB}{
  Since $\oxexpr$ is a value, we can apply \Lemma{lemma:values-dont-change} to
  get $\oxsubtypectx{\oxglobalctx}{\oxemptyctx}{\oxvarctx}{\oxvarctx^\prime}$.
  Applying \Lemma{lemma:stack-validity-related-envs}, then gives us
  $\oxstorevalidity{\oxglobalctx}{\oxvarctx^\prime}{\oxstore}$.

  Then, we can apply \Lemma{lemma:place-exprs-reduce} to
  $\oxmusafety{\oxtvarctx}{\oxvarctx^\prime}{\oxmuta}{\oxplaceexpr}{
    \oxset{\overline{\oxloan}}
  }$, $\oxcomputetynoprov{\oxtvarctx}{\oxvarctx^\prime}{\oxmuta}{\oxplaceexpr}
  {\oxxitype}$, and $\oxstorevalidity{\oxglobalctx}{\oxvarctx^\prime}{\oxstore}$
  to get $\oxnorm{\oxstore}{\oxplaceexpr}{\oxreferent}{\_}{\oxvalue}$. By
  \Lemma{lemma:canonical-forms}, we know that
  $\oxvalue =  \oxarr{\oxvalue_0 \oxdotsc \oxvalue_n}$ since the type tells us
  the shape of the resultant value.

  Since we wish to step with one of \oxname{E-BorrowIndex} and
  \oxname{E-BorrowIndexOOB}, we should observe that we now have
  their shared requirement: $\oxnorm{\oxstore}{\oxplaceexpr}{\oxreferent}{\_}{
    \oxarr{\oxvalue_0 \oxdotsc \oxvalue_n}
  }$. Their other obligations are a bounds check which together are a tautology
  (i.e. one of them must hold). Thus, we can step with the appropriate rule
  based on whether or not the bounds check succeeds.
}

\oxprogresscaseinductive{T-BorrowSlice}{\TBorrowSlice}{
  The proof proceeds along similar lines as for \oxname{T-BorrowIndex}. We
  proceed based on whether or not $\oxexpr_1$ and $\oxexpr_2$ are values.

  If $\oxexpr_1$ is not a value, then we can decompose our whole expression into
  the evaluation context
  $\oxref{\oxprov}{\oxmuta}{\oxslice{\oxplaceexpr}{\oxhole}{\oxexpr_2}}$ and
  redex $\oxexpr_1$. Then, by applying our induction hypothesis to $\oxexpr_1$,
  we know either that $\oxexpr_1$ steps to some $\oxexpr_1^\prime$ or is an
  $\oxkey{abort!}$ expression. In the former case, this satisfies our
  requirement since we can plug $\oxexpr^\prime$ back into our evaluation
  context. In the latter case, we can step with \oxname{E-EvalCtxAbort}.


  If $\oxexpr_1$ is a value and $\oxexpr_2$ is not a value, then we can
  decompose our whole expression into the evaluation context
  $\oxref{\oxprov}{\oxmuta}{\oxslice{\oxplaceexpr}{\oxvalue_1}{\oxhole}}$ and
  redex $\oxexpr_2$. Then, by applying our induction hypothesis to $\oxexpr_2$,
  we know either that $\oxexpr_2$ steps to some $\oxexpr_2^\prime$ or is an
  $\oxkey{abort!}$ expression. In the former case, this satisfies our
  requirement since we can plug $\oxexpr^\prime$ back into our evaluation
  context. In the latter case, we can step with \oxname{E-EvalCtxAbort}.
}{
  If $\oxexpr_1$ and $\oxexpr_2$ are values, we would like to step with one of:
}{\EBorrowSlice \and \EBorrowSliceOOB}{
  Since $\oxexpr_1$ is a value, we can apply \Lemma{lemma:values-dont-change} to
  get $\oxsubtypectx{\oxglobalctx}{\oxemptyctx}{\oxvarctx}{\oxvarctx_1}$.
  Then, since $\oxexpr_2$ is also a value, we can apply
  \Lemma{lemma:values-dont-change} to get
  $\oxsubtypectx{\oxglobalctx}{\oxemptyctx}{\oxvarctx_1}{\oxvarctx_2}$. Then, by
  transitivity, we get
  $\oxsubtypectx{\oxglobalctx}{\oxemptyctx}{\oxvarctx}{\oxvarctx_2}$.
  Then, applying \Lemma{lemma:stack-validity-related-envs} gives us
  $\oxstorevalidity{\oxglobalctx}{\oxvarctx_2}{\oxstore}$.

  Then, we can apply \Lemma{lemma:place-exprs-reduce} to
  $\oxmusafety{\oxtvarctx}{\oxvarctx_2}{\oxmuta}{\oxplaceexpr}{
    \oxset{\overline{\oxloan}}
  }$, $\oxcomputetynoprov{\oxtvarctx}{\oxvarctx_2}{\oxmuta}{\oxplaceexpr}
  {\oxtslice{\oxsitype}}$, and $\oxstorevalidity{\oxglobalctx}{\oxvarctx_2}{
    \oxstore
  }$ to get $\oxnorm{\oxstore}{\oxplaceexpr}{\oxreferent}{\_}{\oxvalue}$. By
  \Lemma{lemma:canonical-forms}, we know that
  $\oxvalue =  \oxarr{\oxvalue_0 \oxdotsc \oxvalue_n}$ since the type tells us
  the shape of the resultant value.

  Since we wish to step with one of \oxname{E-BorrowSlice} and
  \oxname{E-BorrowSliceOOB}, we should observe that we now have
  their shared requirement: $\oxnorm{\oxstore}{\oxplaceexpr}{\oxreferent}{\_}{
    \oxarr{\oxvalue_0 \oxdotsc \oxvalue_n}
  }$. Their other obligations are a bounds check which together are a tautology
  (i.e. one of them must hold). Thus, we can step with the appropriate rule
  based on whether or not the bounds check succeeds.
}

\oxprogresscaseinductive{T-IndexCopy}{\TIndexCopy}{
  We proceed based on whether or not $\oxexpr$ is a value. If it is not, we can
  decompose our expression into the evaluation context
  $\oxindex{\oxplaceexpr}{\oxhole}$ and redex $\oxexpr$. Then, by applying our
  induction hypothesis to $\oxexpr$, we know either that $\oxexpr$ steps to
  some $\oxexpr^\prime$ or is an $\oxkey{abort!}$ expression. In the former
  case, this satisfies our requirement since we can plug $\oxexpr^\prime$ back
  into our evaluation context. In the latter case, we can step with
  \oxname{E-EvalCtxAbort}.
}{If $\oxexpr$ is a value, we would like to step with one of:}{
  \EIndexCopy \and \EIndexCopyOOB
}{
  Since $\oxexpr$ is a value, we can apply \Lemma{lemma:values-dont-change} to
  get $\oxsubtypectx{\oxglobalctx}{\oxemptyctx}{\oxvarctx}{\oxvarctx^\prime}$.
  Applying \Lemma{lemma:stack-validity-related-envs}, then gives us
  $\oxstorevalidity{\oxglobalctx}{\oxvarctx^\prime}{\oxstore}$.

  Then, we can apply \Lemma{lemma:place-exprs-reduce} to
  $\oxmusafety{\oxtvarctx}{\oxvarctx^\prime}{\oxmuta}{\oxplaceexpr}{
    \oxset{\overline{\oxloan}}
  }$, $\oxcomputetynoprov{\oxtvarctx}{\oxvarctx^\prime}{\oxmuta}{\oxplaceexpr}
  {\oxxitype}$, and $\oxstorevalidity{\oxglobalctx}{\oxvarctx^\prime}{\oxstore}$
  to get $\oxnorm{\oxstore}{\oxplaceexpr}{\_}{\_}{\oxvalue}$. By
  \Lemma{lemma:canonical-forms}, we know that
  $\oxvalue =  \oxarr{\oxvalue_0 \oxdotsc \oxvalue_n}$ since the type tells us
  the shape of the resultant value.

  Since we wish to step with one of \oxname{E-IndexCopy} and
  \oxname{E-IndexCopyOOB}, we should observe that we now have
  their shared requirement: $\oxnorm{\oxstore}{\oxplaceexpr}{\_}{\_}{
    \oxarr{\oxvalue_0 \oxdotsc \oxvalue_n}
  }$. Their other obligations are a bounds check which together are a tautology
  (i.e. one of them must hold). Thus, we can step with the appropriate rule
  based on whether or not the bounds check succeeds.
}

\oxprogresscaseinductive{T-Seq}{\TSeq}{
  We proceed based on whether or not $\oxexpr_1$ is a value. If it is not, we can
  decompose our expression into the evaluation context
  $\oxseq{\oxhole}{\oxexpr_2}$ and redex $\oxexpr_1$. Then, by applying our
  induction hypothesis to $\oxexpr_1$, we know either that $\oxexpr_1$ steps to
  some $\oxexpr_1^\prime$ or is an $\oxkey{abort!}$ expression. In the former
  case, this satisfies our requirement since we can plug $\oxexpr_1^\prime$ back
  into our evaluation context. In the latter case, we can step with
  \oxname{E-EvalCtxAbort}.
}{If $\oxexpr_1$ is a value, we can step with:}{\ESeq}{}

\oxprogresscaseinductive{T-Branch}{\TBranch}{
  We proceed based on whether or not $\oxexpr_1$ is a value. If it is not, we can
  decompose our expression into the evaluation context
  $\oxbranch{\oxhole}{\oxexpr_2}{\oxexpr_3}$ and redex $\oxexpr_1$. Then, by
  applying our induction hypothesis to $\oxexpr_1$, we know either that
  $\oxexpr_1$ steps to some $\oxexpr_1^\prime$ or is an $\oxkey{abort!}$
  expression. In the former case, this satisfies our requirement since we can
  plug $\oxexpr_1^\prime$ back into our evaluation context. In the latter case,
  we can step with \oxname{E-EvalCtxAbort}.
}{If $\oxexpr_1$ is a value, we would like to step with one of:}{
  \EIfTrue \and \EIfFalse
}{
  Since $\oxexpr_1$ is a value, applying \Lemma{lemma:canonical-forms} tells us
  that $\oxexpr_1$ is either $\oxtrue$ or $\oxfalse$. In the former case, we can
  step with \oxname{E-IfTrue} and in the latter case, we can step with
  \oxname{E-IfFalse}
}

\oxprogresscaseinductive{T-Let}{\TLet}{
  We proceed based on whether or not $\oxexpr_1$ is a value. If it is not, we can
  decompose our expression into the evaluation context
  $\oxlet{\oxid}{\oxsitype_a}{\oxhole}{\oxexpr_2}$ and redex $\oxexpr_1$. Then,
  by applying our induction hypothesis to $\oxexpr_1$, we know either that
  $\oxexpr_1$ steps to some $\oxexpr_1^\prime$ or is an $\oxkey{abort!}$
  expression. In the former case, this satisfies our requirement since we can
  plug $\oxexpr_1^\prime$ back into our evaluation context. In the latter case,
  we can step with \oxname{E-EvalCtxAbort}.
}{If $\oxexpr_1$ is a value, we can step with:}{\ELet}{}

\oxprogresscaseinductive{T-LetRegion}{\TLetProv}{
  We proceed based on whether or not $\oxexpr$ is a value. If it is not, we can
  decompose our expression into the evaluation context
  $\oxletrgn{\oxcprov}{\oxhole}$ and redex $\oxexpr$. Then, by applying our
  induction hypothesis to $\oxexpr$, we know either that $\oxexpr$ steps to some
  $\oxexpr^\prime$ or is an $\oxkey{abort!}$ expression. In the former case,
  this satisfies our requirement since we can plug $\oxexpr^\prime$ back into
  our evaluation context. In the latter case, we can step with
  \oxname{E-EvalCtxAbort}.
}{If $\oxexpr$ is a value, we can step with:}{\ELetProv}{}

\oxprogresscaseinductive{T-Assign}{\TAssign}{
  We proceed based on whether or not $\oxexpr$ is a value. If it is not, we can
  decompose our expression into the evaluation context
  $\oxassign{\oxplaceexpr}{\oxhole}$ and redex $\oxexpr$. Then, by applying our
  induction hypothesis to $\oxexpr$, we know either that $\oxexpr$ steps to some
  $\oxexpr^\prime$ or is an $\oxkey{abort!}$ expression. In the former case,
  this satisfies our requirement since we can plug $\oxexpr^\prime$ back into
  our evaluation context. In the latter case, we can step with
  \oxname{E-EvalCtxAbort}.
}{
  If $\oxexpr$ is a value, we would like to step with:
}{\EAssign}{
  Since $\oxexpr$ is a value, we can apply \Lemma{lemma:values-dont-change} to
  get $\oxsubtypectx{\oxglobalctx}{\oxemptyctx}{\oxvarctx}{\oxvarctx^\prime}$.
  Applying \Lemma{lemma:stack-validity-related-envs}, then gives us
  $\oxstorevalidity{\oxglobalctx}{\oxvarctx^\prime}{\oxstore}$.

  Then, we can apply \Lemma{lemma:place-exprs-reduce} to
  $\oxmusafety{\oxtvarctx}{\oxvarctx^\prime}{\oxmuta}{\oxplaceexpr}{
    \oxset{\overline{\oxloan}}
  }$, $\oxcomputetynoprov{\oxtvarctx}{\oxvarctx^\prime}{\oxmuta}{\oxplaceexpr}
  {\oxxitype}$, and $\oxstorevalidity{\oxglobalctx}{\oxvarctx^\prime}{\oxstore}$
  to get $\oxnorm{\oxstore}{\oxplaceexpr}{\oxreferent}{\oxvaluectx}{\_}$.
}

\oxprogresscaseinductive{T-ForArray}{\TForArray}{
  We proceed based on whether or not $\oxexpr_1$ is a value. If it is not, we can
  decompose our expression into the evaluation context
  $\oxfor{\oxid}{\oxhole}{\oxexpr_2}$ and redex $\oxexpr_1$. Then, by applying our
  induction hypothesis to $\oxexpr_1$, we know either that $\oxexpr_1$ steps to
  some $\oxexpr_1^\prime$ or is an $\oxkey{abort!}$ expression. In the former
  case, this satisfies our requirement since we can plug $\oxexpr_1^\prime$ back
  into our evaluation context. In the latter case, we can step with
  \oxname{E-EvalCtxAbort}.
}{If $\oxexpr_1$ is a value, we would like to step with one of:}{
  \EForArray \and \EForEmptyArray
}{
  Since $\oxexpr_1$ is a value, then by \Lemma{lemma:canonical-forms}, we know
  that $\oxexpr_1$ is of the form $\oxarr{\oxvalue_1 \oxdotsc \oxvalue_n}$. If
  $n > 0$, then we can step with \oxname{E-ForArray}, and if $n = 0$, then we
  can step with \oxname{E-ForEmptyArray}.
}

\oxprogresscaseinductive{T-ForSlice}{\TForSlice}{
  We proceed based on whether or not $\oxexpr_1$ is a value. If it is not, we can
  decompose our expression into the evaluation context
  $\oxfor{\oxid}{\oxhole}{\oxexpr_2}$ and redex $\oxexpr_1$. Then, by applying our
  induction hypothesis to $\oxexpr_1$, we know either that $\oxexpr_1$ steps to
  some $\oxexpr_1^\prime$ or is an $\oxkey{abort!}$ expression. In the former
  case, this satisfies our requirement since we can plug $\oxexpr_1^\prime$ back
  into our evaluation context. In the latter case, we can step with
  \oxname{E-EvalCtxAbort}.
}{If $\oxexpr_1$ is a value, we would like to step with one of:}{
  \EForSlice \and \EForEmptySlice
}{
  If $\oxexpr_1$ is a value, then by \Lemma{lemma:canonical-forms}, we know that
  $\oxexpr_1$ is of the form $\oxsliceptr{\oxreferent}{\oxnum_1}{\oxnum_2}$.
  Further, by inversion of \oxname{T-Pointer} for the typing derivation of
  $\oxexpr_1$, we get $\oxreferentvalidity{\oxglobalctx}{\oxvarctx}{
    \oxslice{\oxreferent}{i}{j}
  }{\oxtslice{\oxsitype}}$. By inversion of \oxname{WF-RefSliceArray} or
  \oxname{WF-RefSliceSlice} (one of which must apply since the referent ends in
  a slice), we know that $i \leq j$. If $i < j$, we step with
  \oxname{E-ForSlice} and if $i = j$, we step with \oxname{E-ForEmptySlice}.
}

\oxprogresscasewide{T-Closure}{\TClosure}{\EClosure}{
  Since $\oxfvars{\cdot}$ and $\oxfncvars{\oxstore}{\cdot}$ are total, we can
  always step with \oxname{E-Closure}.
}

\oxprogresscaseinductive{T-AppClosure}{\TAppClosure}{
  We proceed based on whether or not $\oxexpr_f$ is a value. If it is not, we can
  decompose our expression into the evaluation context
  $\oxapply{\oxhole}{\oxexpr_1 \oxdotsc \oxexpr_n}$ and redex
  $\oxexpr_f$. Then, by applying our induction hypothesis to $\oxexpr_f$, we
  know either that $\oxexpr_f$ steps to some $\oxexpr_f^\prime$ or is an
  $\oxkey{abort!}$ expression. In the former case, this satisfies our
  requirement since we can plug $\oxexpr_f^\prime$ back into our evaluation
  context. In the latter case, we can step with \oxname{E-EvalCtxAbort}.

  Next, we'll proceed based on whether or not each expression $\oxexpr_i$ is a
  value. If any of them are not, we can decompose our expression into the
  evaluation context $\oxapply{\oxvalue_f}{\oxvalue_1 \oxdotsc \oxvalue_m
    \oxcomma \oxhole \oxcomma \oxexpr_1 \oxdotsc \oxexpr_{n^\prime}}$ and redex
  $\oxexpr_i$. Then, by applying our induction hypothesis to $\oxexpr_i$, we
  know either that $\oxexpr_i$ steps to some $\oxexpr_i^\prime$ or is an
  $\oxkey{abort!}$ expression. In the former case, this satisfies our
  requirement since we can plug $\oxexpr_i^\prime$ back into our evaluation
  context. In the latter case, we can step with \oxname{E-EvalCtxAbort}.
}{
  If $\oxexpr_f$ is a value and every $\oxexpr_i$ is a value, we would like to
  step with one of:
}{\EApp}{
  Since $\oxexpr_f$ is a value, then by \Lemma{lemma:canonical-forms}, we know
  that it has the form
  $\oxclosureval{\oxstore_c}{ \oxascribe{\oxid_1}{\oxsitype_1} \oxdotsc
    \oxascribe{\oxid_n}{\oxsitype_n} }{\oxsitype_r}{\oxexpr}$.
  Then, since all of the $e_i$ are values, then we can step using
  \oxname{E-AppClosure}.
}

\oxprogresscaseinductive{T-AppFunction}{\TApp}{
  We proceed based on whether or not $\oxexpr_f$ is a value. If it is not, we can
  decompose our expression into the evaluation context
  $\oxapp{\oxhole}{\overline{\oxenv} \oxcomma \overline{\oxprov} \oxcomma
    \overline{\oxsitype}}{\oxexpr_1 \oxdotsc \oxexpr_n}$ and redex
  $\oxexpr_f$. Then, by applying our induction hypothesis to $\oxexpr_f$, we
  know either that $\oxexpr_f$ steps to some $\oxexpr_f^\prime$ or is an
  $\oxkey{abort!}$ expression. In the former case, this satisfies our
  requirement since we can plug $\oxexpr_f^\prime$ back into our evaluation
  context. In the latter case, we can step with \oxname{E-EvalCtxAbort}.

  Next, we'll proceed based on whether or not each expression $\oxexpr_i$ is a
  value. If any of them are not, we can decompose our expression into the
  evaluation context
  $\oxapp{\oxvalue_f}{\overline{\oxenv} \oxcomma \overline{\oxprov} \oxcomma \overline{\oxsitype} }{\oxvalue_1 \oxdotsc \oxvalue_m \oxcomma \oxhole \oxcomma \oxexpr_1 \oxdotsc \oxexpr_{n^\prime}}$
  and redex $\oxexpr_i$. Then, by applying our induction hypothesis to
  $\oxexpr_i$, we know either that $\oxexpr_i$ steps to some $\oxexpr_i^\prime$
  or is an $\oxkey{abort!}$ expression. In the former case, this satisfies our
  requirement since we can plug $\oxexpr_i^\prime$ back into our evaluation
  context. In the latter case, we can step with \oxname{E-EvalCtxAbort}. }{ If
  $\oxexpr_f$ is a value and every $\oxexpr_i$ is a value, we would like to step
  with one of: }{\EAppFun}{ Since $\oxexpr_f$ is a value, then by
  \Lemma{lemma:canonical-forms}, we know that it has the form
  $\oxfnname$.
  Then, since all of the $e_i$ are values, then we can step using
  \oxname{E-AppFunction}.
}

\oxprogresscasevalue{T-Unit}{\TUnit}{
  By inspection of the value grammar, we know that $\oxunit$ is already a value.
}

\oxprogresscasevalue{T-u32}{\TuThreeTwo}{
  By inspection of the value grammar, we know that $\oxnum$ is already a value.
}

\oxprogresscasevalue{T-True}{\TTrue}{
  By inspection of the value grammar, we know that $\oxtrue$ is already a value.
}

\oxprogresscasevalue{T-False}{\TFalse}{
  By inspection of the value grammar, we know that $\oxfalse$ is already a
  value.
}

\oxprogresscasevalue{T-Tuple}{\TTuple}{
  We'll proceed based on whether or not each expression $\oxexpr_i$ is a
  value. If any of them are not, we can decompose our expression into the
  evaluation context $\oxprod{\oxvalue_1 \oxdotsc \oxvalue_m \oxcomma \oxhole
    \oxcomma \oxexpr_1 \oxdotsc \oxexpr_{n^\prime}}$ and redex
  $\oxexpr_i$. Then, by applying our induction hypothesis to $\oxexpr_i$, we
  know either that $\oxexpr_i$ steps to some $\oxexpr_i^\prime$ or to an
  $\oxkey{abort!}$ expression. In either case, this satisfies our requirement,
  since we can plug $\oxexpr_i^\prime$ back into our evaluation context.

  If every expression $e_i$ is a value, then the whole expression is a value by
  the definition of values.
}

\oxprogresscasevalue{T-Array}{\TArray}{
  We'll proceed based on whether or not each expression $\oxexpr_i$ is a
  value. If any of them are not, we can decompose our expression into the
  evaluation context $\oxarr{\oxvalue_1 \oxdotsc \oxvalue_m \oxcomma \oxhole
    \oxcomma \oxexpr_1 \oxdotsc \oxexpr_{n^\prime}}$ and redex
  $\oxexpr_i$. Then, by applying our induction hypothesis to $\oxexpr_i$, we
  know either that $\oxexpr_i$ steps to some $\oxexpr_i^\prime$ or to an
  $\oxkey{abort!}$ expression. In either case, this satisfies our requirement,
  since we can plug $\oxexpr_i^\prime$ back into our evaluation context.

  If every expression $e_i$ is a value, then the whole expression is a value by
  the definition of values.
}

\oxprogresscasevalue{T-Abort}{\TAbort}{
  By definition, $\oxabort{\oxdots}$ is an $\oxkey{abort!}$ expression.
}

% \oxprogresscasevalue{T-TupleStruct}{\TTupleStruct}{
%   We'll proceed based on whether or not each expression $\oxexpr_i$ is a
%   value. If any of them are not, we can decompose our expression into the
%   evaluation context $\oxtupstruct{
%     \oxturbofish{\oxdatatype}{\overline{\oxprov} \oxcomma \overline{\oxsitype}}
%   }{\oxvalue_1 \oxdotsc \oxvalue_m \oxcomma \oxhole
%     \oxcomma \oxexpr_1 \oxdotsc \oxexpr_{n^\prime}}$ and redex
%   $\oxexpr_i$. Then, by applying our induction hypothesis to $\oxexpr_i$, we
%   know either that $\oxexpr_i$ steps to some $\oxexpr_i^\prime$ or to an
%   $\oxkey{abort!}$ expression. In either case, this satisfies our requirement,
%   since we can plug $\oxexpr_i^\prime$ back into our evaluation context.

%   If every expression $e_i$ is a value, then the whole expression is a value by
%   the definition of values.
% }

% \oxprogresscasevalue{T-RecordStruct}{\TRecordStruct}{
%   We'll proceed based on whether or not each expression $\oxexpr_i$ is a
%   value. If any of them are not, we can decompose our expression into the
%   following evaluation context:
%   \begin{center}
%     $\oxrecstruct{
%       \oxturbofish{\oxdatatype}{\overline{\oxprov} \oxcomma \overline{\oxsitype}}
%     }{\oxid_1 \oxcolon \; \oxvalue_1 \oxdotsc \oxid_m \oxcolon \; \oxvalue_m \oxcomma
%       \oxid_{m+1} \oxcolon \; \oxhole
%       \oxcomma \oxid_{m+2} \oxcolon \; \oxexpr_1 \oxdotsc \oxid_n \oxcolon \;
%       \oxexpr_{n^\prime}}$
%   \end{center}
%   and the redex $\oxexpr_i$. Then, by applying our induction hypothesis to
%   $\oxexpr_i$, we know either that $\oxexpr_i$ steps to some $\oxexpr_i^\prime$
%   or to an $\oxkey{abort!}$ expression. In either case, this satisfies our
%   requirement, since we can plug $\oxexpr_i^\prime$ back into our evaluation
%   context.

%   If every expression $e_i$ is a value, then the whole expression is a value by
%   the definition of values.
% }

\oxprogresscaseinductive{T-Framed}{\TFramed}{
  We proceed based on whether or not $\oxexpr$ is a value. If it is not, we can
  decompose our expression into the evaluation context
  $\oxframed{\oxhole}$ and redex $\oxexpr$. Then, by applying our
  induction hypothesis to $\oxexpr$, we know either that $\oxexpr$ steps to some
  $\oxexpr^\prime$ or to an $\oxkey{abort!}$ expression. In either case, this
  satisfies our requirement, since we can plug $\oxexpr^\prime$ back into our
  evaluation context.
}{If $\oxexpr$ is a value, then we would like to step with:}{\EFramed}{
  In order to do so, we need to know $\oxid \in \oxdomain{\oxstore}$.
  Fortunately, we know from our assumption that
  \oxstorevalidity{\oxglobalctx}{\oxvarctx}{\oxstore} (via
  \oxname{WF-Stack}). The premise of \oxname{WF-Stack} tells us that
  $\oxdomain{\oxstore} = \oxdomain{\oxvarctx}$, and thus the $\oxid \in
  \oxdomain{\oxvarctx}$ from the premise of \oxname{T-Framed} is sufficient to
  tell us that $\oxid \in \oxdomain{\oxstore}$. Thus, we can step with
  \oxname{E-Framed}.
}

\oxprogresscasevalue{T-Pointer}{\TPointer}{
  By inspection of the value grammar, we know that
  $\oxptr{\oxplace}$ is already a value.
}

\oxprogresscasevalue{T-ClosureValue}{\TClosureVal}{
  By inspection of the value grammar, we know that
  $\oxclosureval{\oxstore}{ \oxascribe{\oxid_1}{\oxsitype_1} \oxdotsc
    \oxascribe{\oxid_n}{\oxsitype_n} }{\oxsitype_r}{\oxexpr}$ is already a
  value.
}

\oxprogresscasevalue{T-Dead}{\TUninit}{
  The type $\oxsidtype$ is not in the grammar of $\oxsitype$. Thus, we have a
  contradiction.
}

\oxprogresscasevalue{T-Drop}{\TDrop}{
  By \oxname{R-Env}, we have that $\oxsubtypectx{\oxglobalctx}{\oxemptyctx}
  {\oxvarctx}{\oxtupdate{\oxvarctx}{\oxplace}{\oxsidtype_\oxplace}}$. Then,
  applying \Lemma{lemma:stack-validity-related-envs} with
  $\oxstorevalidity{\oxglobalctx}{\oxvarctx}{\oxstore}$ (from our premise) gives
  us $\oxstorevalidity{\oxglobalctx}{
    \oxtupdate{\oxvarctx}{\oxplace}{\oxsidtype_\oxplace}
  }{\oxstore}$. We can then apply our induction hypothesis to this and
  $\oxtypjudge{\oxglobalctx}{\oxemptyctx}{\oxkontctx}{
    \oxtupdate{\oxvarctx}{\oxplace}{\oxsidtype_\oxplace}
  }{\oxexpr}{\oxsxtype}{\oxvarctx_f}$ to reach our goal.
}

\oxprogresscasevalue{T-Left}{\TInl}{
  We proceed based on whether or not $\oxexpr$ is a value. If it is not, we can
  decompose our expression into the evaluation context
  $\oxinl{\oxsitype_1}{\oxsitype_2}{\oxhole}$ and redex $\oxexpr$. Then, by
  applying our induction hypothesis to $\oxexpr$, we know either that $\oxexpr$
  steps to some $\oxexpr^\prime$ or to an $\oxkey{abort!}$ expression. In either
  case, this satisfies our requirement, since we can plug $\oxexpr^\prime$ back
  into our evaluation context.

  If $\oxexpr$ is a value, then we know that the whole expression
  $\oxinl{\oxsitype_1}{\oxsitype_2}{\oxvalue}$ is a value and thus we are done.
}

\oxprogresscasevalue{T-Right}{\TInr}{
  We proceed based on whether or not $\oxexpr$ is a value. If it is not, we can
  decompose our expression into the evaluation context
  $\oxinr{\oxsitype_1}{\oxsitype_2}{\oxhole}$ and redex $\oxexpr$. Then, by
  applying our induction hypothesis to $\oxexpr$, we know either that $\oxexpr$
  steps to some $\oxexpr^\prime$ or to an $\oxkey{abort!}$ expression. In either
  case, this satisfies our requirement, since we can plug $\oxexpr^\prime$ back
  into our evaluation context.

  If $\oxexpr$ is a value, then we know that the whole expression
  $\oxinr{\oxsitype_1}{\oxsitype_2}{\oxvalue}$ is a value and thus we are done.
}

\oxprogresscaseinductive{T-Match}{\TMatch}{
  We proceed based on whether or not $\oxexpr$ is a value. If it is not, we can
  decompose our expression into the evaluation context
  $\oxmatch{\oxhole}{\oxid_1}{\oxexpr_1}{\oxid_2}{\oxexpr_2}$ and redex
  $\oxexpr$. Then, by applying our induction hypothesis to $\oxexpr$, we know
  either that $\oxexpr$ steps to some $\oxexpr^\prime$ or to an $\oxkey{abort!}$
  expression. In either case, this satisfies our requirement, since we can plug
  $\oxexpr^\prime$ back into our evaluation context.
}{If $\oxexpr$ is a value, then we would like to step with either:}{
  \EMatchLeft \and \EMatchRight
}{
  Since $\oxexpr$ is a value, applying \Lemma{lemma:canonical-forms} tells us
  that $\oxexpr$ is either of the form $\oxinl{\oxsitype_1}{\oxsitype_2}{\oxvalue}$ or
  $\oxinr{\oxsitype_1}{\oxsitype_2}{\oxvalue}$. In the former case, we can
  step with \oxname{E-MatchLeft} and in the latter case, we can step with
  \oxname{E-MatchRight}.
}

%%% Local Variables:
%%% mode: latex
%%% fill-column: 80
%%% End:
